\section{Local vertex conventions and the crossed prefactor}
\label{app:vertex}
We distinguish the vertex reconstructed from the mode integral from the
literal expression in (A.7) of~\cite{KL}. This appendix fixes the convention
used in \eqref{eq:local-mode}; it is not an author-confirmed erratum.
Set $r_j=s_j\tanh\pi s_j$, $c_j=\cosh\pi s_j$, and define
\begin{align}
 G_4(s_1,s_2,s_3)
 &=\Gamma\!\left(\tfrac14+\tfrac i2(-s_1+s_2+s_3)\right)
   \Gamma\!\left(\tfrac14+\tfrac i2(-s_1-s_2+s_3)\right)\notag\\
 &\quad\times
   \Gamma\!\left(\tfrac14+\tfrac i2(-s_1+s_2-s_3)\right)
   \Gamma\!\left(\tfrac14+\tfrac i2(s_1+s_2+s_3)\right),\\
 D(s_1,s_2,s_3)
 &=\Gamma(1/2-is_1)\Gamma(1/2+is_2)\Gamma(1/2+is_3).
\end{align}
Then the reconstructed vertex is
\begin{equation}
 \mathcal V^{s_1}_{s_2,s_3,\mathrm{rec}}
 =e^{-i\pi/4}\frac{c_1\sqrt{r_1r_2r_3}}
                    {8\pi\sqrt{c_1c_2c_3}}\frac{G_4}{D}
 =\frac{\mathcal V^{s_1}_{s_2,s_3,\mathrm{literal}}}
        {\Gamma(1/2+is_3)^2}.
 \label{eq:reconstructed-vertex}
\end{equation}
The second equality states the discrepancy explicitly.

Here is its analytic normalization check. Write
\[
 P_j(z)=P_{-1/2+is_j}(z),\qquad
 \Psi_0^{s_j}(t)=\frac{\sqrt{r_j}\,P_j(i\sinh t)}{2\sqrt{\pi c_j}}.
\]
Closing the vertical contour to the right and using the Legendre jump
relation gives
\begin{equation}
 \int_{\mathbb R}\cosh t\,P_1(-i\sinh t)P_2(i\sinh t)P_3(i\sinh t)\dd t
 =2c_1\int_1^\infty P_1(u)P_2(u)P_3(u)\dd u.
\end{equation}
The endpoint circle is $O(\epsilon|\log\epsilon|)$ and the large
semicircle is $O(R^{-1/2})$, for fixed positive labels. The triple-Legendre
identity of~\cite{Bros} is
\begin{equation}
 \int_1^\infty P_1P_2P_3\dd u
 =\frac{\Gamma_8\,c_1c_2c_3}{4\pi^5},\qquad
 \Gamma_8=\prod_{\epsilon_j=\pm1}
 \Gamma\!\left(\tfrac14+\tfrac i2\sum_{j=1}^3\epsilon_js_j\right).
\end{equation}
Including the angular integral yields
\begin{equation}
 T_{000}=\frac{c_1\sqrt{r_1r_2r_3}\sqrt{c_1c_2c_3}}
                     {8\pi^{11/2}}\Gamma_8.
\end{equation}
The circle three-point integral in the same convention gives
$C_{S,000}^{\mathrm{phys}}=e^{-i\pi/4}G_4/(\pi^{3/2}D)$.
Dividing by its conjugate proves \eqref{eq:reconstructed-vertex} at zero
magnetic labels. It fixes all finite-mode vertices as follows: the complete
principal three-leg invariant has $S$ and $A$ channels; the position integral
is even under simultaneous angular reflection, while the explicit
$A$-channel sign product is odd. The $A$ coefficient vanishes and the
nonzero $S,000$ coefficient fixes the remaining scalar. This uses the
stated reflection and half-phase conventions; the wording of footnote 33
in~\cite{KL} should not be used to override its explicit antisymmetric
kernel. Continuity on compact positive label sets extends the
almost-everywhere decomposition identity to fixed positive labels.

For the crossing graph put $(s_1,s_2,s_3,s_4)=(d,a,b,c)$ and let
$A=\nu_A$, $B=\nu_B$. The four scalar factors in (B.2) of~\cite{KL} are
\begin{equation}
 P_{\mathcal V}=
 \overline{\mathcal V^d_{A,a,\mathrm{rec}}}\,
 \overline{\mathcal V^a_{B,b,\mathrm{rec}}}\,
 \mathcal V^c_{A,b,\mathrm{rec}}\,
 \mathcal V^d_{B,c,\mathrm{rec}}.
\end{equation}
Their half phases cancel. Define the sixteen-gamma numerator explicitly by
\begin{equation}
 \Gamma_{16}^{B}
 =\overline{G_4(d,A,a)}\,\overline{G_4(a,B,b)}
                    G_4(c,A,b)G_4(d,B,c).
\end{equation}
The six pairs of denominator gamma functions give
$\prod_{s=d,a,b,c,A,B}c_s/\pi^6$. Thus
\begin{equation}
 P_{\mathcal V}
 =\frac{r_dr_ar_br_cr_Ar_B\,c_d^2c_ac_c}{4096\pi^{10}}
                         \Gamma_{16}^{B}.
 \label{eq:four-vertex-product}
\end{equation}
The numerator is exactly the collection of arguments in (B.27).

Let $S(d,a,b,c,A,B)$ denote the raw recoupling function in (B.17),
with its stated conjugation orientation. In normalized trees the $S/S$
subkernel acquires $\sqrt{w(a)w(c)}$. It is a subkernel of the complete
unitary Racah transform, not by itself a unitary transform.
The scalar local factors contribute the inverse
$w(d)\sqrt{w(a)w(c)}$, together with $\pi^2/(r_Ar_B)$.
Using the physical measure $\rhoP(d)$ and the Hartle--Hawking weight
$\sech^2(\pi d)$, \eqref{eq:four-vertex-product} gives
\begin{equation}
 \frac{r_dr_ar_br_c\,c_ac_c}{4096\pi^{11}}
              \Gamma_{16}^{B}S(d,a,b,c,A,B).
 \label{eq:HH-prefactor}
\end{equation}
This is (B.25)--(B.27), with $a,c$ carrying $\sinh$ and $b,d$
carrying $\tanh$. The calculation uses normalized tree inner products,
not an uncancelled $\delta(0)$ or group volume.

Conjugation orientation can be checked independently. The original modes
are fixed by the antiunitary map
$f(t,\theta)\mapsto\overline{f(-t,-\theta)}$. Applying it to the
absolutely convergent integral \eqref{eq:local-mode} shows that each
zero-time finite-mode local matrix element is real. The complete
zero-time contracted spectral kernel is therefore real as well, by
normalized spectral limits. Reversing bra and ket conjugates both the four
vertex factors and the Racah kernel and gives the same zero-time result.
It does not justify conjugating a nonzero-time phase independently.

For equal matter labels, the assembled symmetric function
$\mathcal S^\nu$ in (A.14) can be cyclically relabeled to $(a,b,c,d)$.
The raw $S$ or its gamma factor alone cannot be relabeled this way.
This distinction is necessary when comparing the local construction with
the archived implementation of the crossed kernel.
